\chapter{Experimental Setup}
\label{chap:setup}

\section{Pipeline Implementation}

\subsection{Unsupervised Anomaly Detection with PatchCore (P3)}

This section reports the concrete technical settings of the method derived in
Chapter~\ref{chap:method}. Each monitored exit is handled by an independent
memory bank built from a single reference frame, reflecting the deployment
reality of a fixed camera with a stable normal scene.

\subsubsection{Backbone and feature extraction}

The backbone is a WideResNet50-2 \cite{Zagoruyko2016} pre-trained on ImageNet and
used frozen (no fine-tuning). Every image is resized to $224 \times 224$ and
normalised with ImageNet statistics. Features are read from blocks $\ell_2$ and
$\ell_3$:

\begin{table}[h]
  \centering
  \begin{tabular}{lcc}
    \hline
    Block & Spatial resolution & Channels \\
    \hline
    $\ell_2$ & $28 \times 28$ & $512$ \\
    $\ell_3$ & $14 \times 14$ & $1024$ \\
    \hline
  \end{tabular}
  \caption{Feature maps read from the frozen backbone at $224\times224$ input.}
  \label{tab:p3-layers}
\end{table}

The $\ell_3$ map is upsampled bilinearly to $28 \times 28$ and concatenated with
$\ell_2$ along the channel axis, giving a tensor of shape
$[\,1536, 28, 28\,]$. The $3 \times 3$ neighbourhood average pooling of
equation~\eqref{eq:nbagg} is applied with unit stride and same-padding, and the
result is flattened to $N_p = 784$ patch vectors of dimension $1536$ per image.
The shallowest block ($56\times56$, low-level texture) and the deepest block
($7\times7$, class-biased abstraction) are not used.

\subsubsection{Build and calibration settings}

The bank is built from $n_{\mathrm{aug}} = 15$ photometric variants of
$I_{\mathrm{ref}}$ (random seed $42$); the original reference is excluded. The
photometric augmentation ranges are listed in Table~\ref{tab:p3-photo}.

\begin{table}[h]
  \centering
  \begin{tabular}{lll}
    \hline
    Transform & Range & Rationale \\
    \hline
    Brightness & $\times[0.85, 1.15]$ & Morning / evening / artificial light \\
    Contrast & $\times[0.90, 1.10]$ & Direct vs.\ diffuse illumination \\
    Saturation & $\times[0.90, 1.10]$ & LED / fluorescent / incandescent \\
    Red-channel shift & $[-8, +8]$ & Colour temperature \\
    Blue-channel shift & $[-20, +20]$ & Colour temperature \\
    Gaussian noise & $\sigma \in [1, 4]$ & Sensor noise \\
    Gaussian blur & $r \in [0.5, 1.5]$, $p=0.3$ & Slight lens defocus \\
    \hline
  \end{tabular}
  \caption{Photometric augmentation applied at build time. No large geometric
  transform is used, as the camera is fixed.}
  \label{tab:p3-photo}
\end{table}

In addition to the photometric transforms, an optional small geometric
perturbation is available for cameras subject to micro-vibration, disabled by
default ($\delta_{\max} = 0$, $\rho_{\max} = 0$). When enabled, every variant
receives a per-variant random translation drawn uniformly in
$[-\delta_{\max}, +\delta_{\max}]$ pixels on each axis and a random rotation drawn
uniformly in $[-\rho_{\max}, +\rho_{\max}]$ degrees, applied by bilinear affine
warping at the original resolution; typical settings are $\delta_{\max} = 8$ px
and $\rho_{\max} = 2^{\circ}$. Unlike the forbidden flips and wide rotations, this
jitter perturbs the whole scene rigidly and teaches the memory to tolerate a small
global displacement at no inference-time cost. It is left off for all experiments
reported in Chapter~\ref{chap:analysis}, where the data-collection protocol (fixed
foot position, braced camera) keeps hand-held variation within the native
tolerance of the neighbourhood pooling and the anomaly-map smoothing.

The raw memory of $n_{\mathrm{aug}} \cdot N_p = 15 \times 784 = 11760$ vectors is
reduced by greedy coreset selection to a fraction $p = 0.01$ (retaining at least
$50$ points), leaving about $117$ vectors. For efficiency the greedy selection
of equation~\eqref{eq:coreset} runs on a sparse random projection to $128$
dimensions (a Johnson-Lindenstrauss embedding), while the final coreset is taken
from the original $1536$-dimensional features.

The threshold is calibrated on $n_{\mathrm{cal}} = 10$ further variants generated
with seed $1000$, disjoint from the bank seed to guarantee statistical
independence; later experiments raise this to $n_{\mathrm{cal}} = 30$ to reduce
the sampling error on $\sigma$. The re-weighting of equation~\eqref{eq:reweight}
uses a neighbourhood of size $b = 9$, and the anomaly map is smoothed with a
Gaussian filter of $\sigma = 4$. The memory bank is serialised to a single file
holding the coreset features together with the calibration parameters $\mu$ and
$\sigma$.

\subsubsection{Anti-leakage separation}

The bank, calibration, and test sets are strictly disjoint
(Table~\ref{tab:p3-leakage}): no data used to build the memory or calibrate the
threshold enters the test phase, and $I_{\mathrm{ref}}$ is reserved as the only
genuinely unseen normal sample.

\begin{table}[h]
  \centering
  \begin{tabular}{llcl}
    \hline
    Set & Seed & Contains $I_{\mathrm{ref}}$? & Role \\
    \hline
    Bank variants ($n_{\mathrm{aug}}=15$) & $42$ & No & Builds the feature memory \\
    Calibration variants ($n_{\mathrm{cal}}=10$) & $1000$ & No & Estimates $\mu$ and $\sigma$ \\
    Normal test & -- & Yes & Independent normal sample \\
    Obstructed test & -- & -- & Anomalous query \\
    \hline
  \end{tabular}
  \caption{Disjoint partition of data across build, calibration, and test.}
  \label{tab:p3-leakage}
\end{table}

\subsubsection{Shadow and light augmentation}

Both disturbance families implement three modalities, summarised in
Table~\ref{tab:p3-disturb}. They are applied after the photometric transforms,
stacking $k_s \in [1,3]$ instances per affected variant. The light family adds
two axes absent from shadows: a diagonal (solar) stripe parameterised by a
continuous incidence angle, and a warm colour shift (red up, blue down) on the
elliptical and diagonal-stripe lights, modelling the chromaticity of direct
sunlight. As a consequence the light subspace is intrinsically richer, so the
raw values of light and shadow false-positive rates are not directly comparable;
only within-category comparisons (across configurations) are legitimate.

\begin{table}[h]
  \centering
  \begin{tabular}{lll}
    \hline
    Modality & Shadow & Light \\
    \hline
    Elliptical & dark ellipse, blurred border & bright ellipse (beam, specular) \\
    Directional & dark edge gradient & bright edge gradient (lateral window) \\
    Stripe & dark band (2 orientations) & bright band (3 orient., solar diagonal) \\
    Colour shift & none & warm tint on elliptical / diagonal \\
    \hline
  \end{tabular}
  \caption{Synthetic disturbance modalities. Build probabilities
  $p_{\mathrm{sh}}, p_{\mathrm{li}} = 0.4$; calibration probabilities
  $p_{\mathrm{sh}}^{\mathrm{cal}}, p_{\mathrm{li}}^{\mathrm{cal}} = 0.1$.}
  \label{tab:p3-disturb}
\end{table}

To formally separate the test-time disturbances from the build-time ones, the
build augmentation draws its shadows and lights from a dedicated random generator
seeded disjointly from the generator that produces the test disturbances, so the
two sets never coincide.

\subsubsection{ROI annotation and connected-component gating}

The frame mask of equation~\eqref{eq:gating} is produced by a three-stage tool
chain. A lightweight interactive annotator draws one polygon per frame element
(left jamb, right jamb, central mullion, threshold, architrave, or a catch-all
label) and writes one JSON file per image, storing polygons in original-image
coordinates to avoid precision loss after resizing. An importer loads these
polygons into the database, matching by relative file path. At evaluation time
the polygon set of the reference frame is rasterised to the original resolution,
the smoothed anomaly map is binarised at $\theta_i$, its $8$-connected components
are extracted, and a component is kept only if its overlap with the frame mask
reaches the minimum $\tau = 10$ pixels. The score is the maximum of the
continuous map over the surviving components, or zero if none survive.

The binarisation threshold used to form the components is the same per-camera
$\theta_i$ used for the decision, so the sensitivity $k$ acts at two stages -
which pixels survive binarisation, and which final scores clear the decision. A
legacy rectangular ROI can be composed in front of the polygon gating: the
rectangle zeroes the map outside the door bounding box first (suppressing floor
reflections and solar shadows), and the polygon rule then filters the remaining
glass-panel variation. The full image is always included in the memory bank:
excluding the glass region would leave its patches without any neighbour at test
time, inflating their distance by mere absence of reference and making the panel
anomalous by construction.

\subsubsection{Pooled-sigma evaluation}

The pooled threshold of equation~\eqref{eq:pooled} is computed in two passes over
the cameras. The first pass builds each bank, calibrates $\mu_i$ and $\sigma_i$,
and scores all test samples without yet deciding their label. The second pass
collects the $\sigma_i$, sets $\sigma_{\mathrm{pop}}$ to their median, and
recomputes for every sample the threshold $\theta_i = \mu_i + k\,\sigma_{\mathrm{pop}}$,
the normalised score $(s - \mu_i)/\sigma_{\mathrm{pop}}$, and the binary decision.
The normalisation makes scores comparable across cameras of different absolute
scale, so that the aggregated AUROC measures intrinsic separability rather than
scale differences. The mode (per-camera or pooled) and the value of
$\sigma_{\mathrm{pop}}$ are recorded with the experiment for reproducibility.
Pooling affects neither inference time nor memory size, since calibration
variants never enter the coreset; its only cost is a one-off computation at build
time.

\section{Application}

[General description of the application, its modules, and overall architecture -
to be completed.]
